\section{Arsitektur Scanner: Buffer, Sentinel, dan Lookahead}
\label{sec:ch07-buffer}

\subsection{Mengapa Buffer Diperlukan}

Lexer membaca program sumber karakter demi karakter. Membaca satu
karakter dari disk pada setiap langkah sangat lambat. Karena itu,
lexer biasanya membaca blok besar ke memori, lalu memprosesnya dari
buffer. Pendekatan ini mengurangi biaya I/O dan memungkinkan
lookahead beberapa karakter tanpa akses disk.

\subsection{Skema Dua Buffer}

Skema klasik menggunakan dua buffer setengah ukuran. Buffer pertama
diisi, lalu buffer kedua, dan seterusnya. Di akhir setiap buffer
ditempatkan sentinel \code{EOF}. Ketika pointer mencapai sentinel,
lexer mengisi buffer berikutnya dan mengatur ulang pointer.

\begin{figure}[H]
\centering
\begin{tikzpicture}[node distance=0mm]
  \node[memorycell, minimum width=4.2cm, minimum height=0.9cm,
        fill=ubbblue!10] (b1) {Buffer 1};
  \node[memorycell, minimum width=1.4cm, minimum height=0.9cm,
        right=of b1, fill=gray!15] (e1) {EOF};
  \node[memorycell, minimum width=4.2cm, minimum height=0.9cm,
        right=of e1, fill=ubbblue!10] (b2) {Buffer 2};
  \node[memorycell, minimum width=1.4cm, minimum height=0.9cm,
        right=of b2, fill=gray!15] (e2) {EOF};
  \draw[decorate, decoration={brace, amplitude=4pt}, thick]
        ([yshift=4mm]b1.north west) -- ([yshift=4mm]b1.north east)
        node[midway, above=2mm, font=\small] {lexeme begin};
  \draw[decorate, decoration={brace, amplitude=4pt}, thick]
        ([yshift=4pt]b2.north west) -- ([yshift=4pt]b2.north east)
        node[midway, above=2mm, font=\small] {forward};
\end{tikzpicture}
\caption{Skema dua buffer dengan sentinel di akhir setiap buffer}
\label{fig:ch07-buffer}
\end{figure}

\subsection{Lexeme Begin dan Forward}

Dua pointer digunakan:
\begin{itemize}[leftmargin=*]
  \item \code{lexeme\_begin} menandai awal lexeme yang sedang dibaca;
  \item \code{forward} bergerak ke depan selama karakter masih cocok
        dengan pattern.
\end{itemize}
Ketika \code{forward} mencapai karakter yang tidak cocok, lexer
menghitung panjang lexeme, membuat token, lalu memindahkan
\code{lexeme\_begin} ke posisi \code{forward}. Jika pattern tidak
cocok sama sekali, lexer menghasilkan token kesalahan dan memajukan
\code{forward} satu karakter.

\subsection{Lookahead}

Lookahead adalah kemampuan melihat karakter berikutnya tanpa
mengonsumsinya. Banyak token memerlukan lookahead satu karakter,
misalnya \code{<} versus \code{<=}. Lookahead yang lebih panjang
diperlukan untuk kasus seperti komentar blok dan string escape.
Implementasi sederhana menyediakan \code{peek()} dan
\code{peek\_next()}.

\begin{infobox}
Pada proyek Subset-C, seluruh berkas dibaca ke memori. Ini
menyederhanakan implementasi dan cukup untuk ukuran program
perkuliahan. Untuk compiler produksi, skema dua buffer atau
\eng{memory-mapped file} lebih hemat memori.
\end{infobox}

\subsection{Kompleksitas}

Dengan buffer, setiap karakter diproses paling banyak dua kali:
sekali saat dibaca ke buffer dan sekali saat dipindai. Kompleksitas
waktu lexer adalah \(O(n)\), dengan \(n\) panjang masukan.
Kompleksitas ruang adalah \(O(b)\), dengan \(b\) ukuran buffer,
ditambah tabel transisi DFA.
